\documentclass[10pt,twocolumn]{article}

\usepackage{fontspec}
\usepackage[letterpaper,
            top=0.75in, bottom=1in,
            left=0.625in, right=0.625in,
            columnsep=0.25in]{geometry}
\usepackage{amsmath,amssymb}
\usepackage{graphicx}
\usepackage{microtype}
\usepackage{fancyhdr}
\usepackage{titlesec}
\usepackage{hyperref}

% ── IEEE-style section headings ──────────────────────────────────────────────
\titleformat{\section}
  {\normalfont\bfseries\scshape\centering}
  {\Roman{section}.}{0.5em}{}

\titleformat{\subsection}
  {\normalfont\itshape}
  {\Alph{subsection}.}{0.5em}{}

\titlespacing{\section}{0pt}{8pt}{4pt}
\titlespacing{\subsection}{0pt}{6pt}{3pt}

% ── Header / Footer ──────────────────────────────────────────────────────────
\pagestyle{fancy}
\fancyhf{}
\renewcommand{\headrulewidth}{0pt}
\fancyfoot[C]{\thepage}

% ── Title block ──────────────────────────────────────────────────────────────
\makeatletter
\renewcommand{\maketitle}{%
  \twocolumn[%
    \begin{@twocolumnfalse}
      \centering
      {\LARGE \bfseries \@title \par}\vspace{0.5em}
      {\normalsize \@author \par}\vspace{0.3em}
      {\small \@date \par}\vspace{1.2em}
    \end{@twocolumnfalse}
  ]%
}
\makeatother

\title{AI -Driven Multi -Objective Optimization for Sustainable Construction Cost and Carbon Reduction}
\author{Author Name \quad \textit{Affiliation} \quad \texttt{email@example.com}}
\date{}

\begin{document}

\maketitle

\begin{abstract}
\boldmath
The construction industry is responsible for nearly 39\% of global energy -related carbon emissions,
with embodied carbon in building materials contributing significantly to lifecycle environmental 
impact. Despite the increasing emphasis on sustainable devel opment and carbon -neutral 
construction, project planning processes continue to prioritize economic efficiency over 
environmental performance. This research presents an AI -driven multi -objective optimization 
framework designed to simultaneously minimize con struction cost and embodied carbon emissions 
using structured Bill of Quantities (BOQ) datasets.  
 
The proposed framework transforms unstructured BOQ documents into machine -readable datasets 
through automated parsing, material classification, and emission factor mapping. Each material entry 
is linked to standardized carbon emission factors derived from established lifecycle assessment (LCA) 
databases. Embodied carbon is computed at item, category, and project levels using quantity -based 
carbon intensity calculations. Cost modeling incorporates real -time pricing structures, enabling dual -
objective evaluati on. 
 
To address the trade -off between economic and environmental objectives, the framework employs 
the Non -dominated Sorting Genetic Algorithm II (NSGA -II), an evolutionary multi -objective 
optimization technique capable of generating Pareto -optimal solution set s. Decision variables include 
material substitutions, quantity adjustments, and alternative construction techniques. The 
optimization process yields a diverse frontier of solutions that provide stakeholders with multiple 
trade -off scenarios rather than a s ingle deterministic output.  
 
Experimental validation was conducted using real -world BOQ datasets from mid -scale residential 
construction projects. Results indicate potential embodied carbon reductions of 18 –32\% with cost 
variations maintained within ±5\% of baseline estimates. The syst em successfully identifies low -
carbon material alternatives such as fly -ash blended cement, recycled steel, and optimized structural 
quantities without compromising project feasibility.  
 
The findings demonstrate that structured BOQ analytics combined with AI -based optimization can 
transform conventional cost estimation documents into intelligent sustainability decision -support 
systems. The proposed framework offers scalability for integrati on into Building Information 
Modeling (BIM) environments and supports policy -driven carbon budgeting strategies in modern 
construction management.
\end{abstract}

\textbf{Index Terms---}Artificial Intelligence, Multi -Objective Optimization, Sustainability, Carbon Emissions,
Construction Management, NSGA -II, BOQ Analysis

\section{Introduction}
The construction industry is one of the largest contributors to global greenhouse gas (GHG)
emissions and resource depletion. According to international environmental assessments, buildings 
and infrastructure collectively account for a substantial share of  energy consumption and nearly 40\% 
of global carbon emissions when both operational and embodied carbon are considered. While 
operational emissions related to heating, cooling, and electricity use have been widely studied and 
optimized over the past decade s, embodied carbon — emissions resulting from material extraction, 
manufacturing, transportation, and construction processes — remains insufficiently addressed 
during early -stage planning.  
 
Embodied carbon is particularly significant in materials such as Portland cement, structural steel, 
bricks, aluminum, and glass. Cement production alone contributes approximately 7 –8\% of global CO₂ 
emissions due to energy -intensive clinker manufacturing an d chemical decomposition processes. 
Similarly, steel production involves high -temperature blast furnace operations, resulting in 
substantial fossil fuel consumption. As construction demand continues to grow, especially in 
developing economies, the cumulati ve impact of embodied emissions presents a major challenge to 
achieving global net -zero targets.  
 
Traditionally, construction planning and cost estimation rely heavily on economic metrics. Engineers 
and project managers prioritize cost optimization, structural safety, regulatory compliance, and 
schedule efficiency. Environmental considerations are ofte n introduced only during later design 
stages or through external certification frameworks such as green building standards. This sequential 
approach limits the potential for impactful carbon reductions, as material choices and structural 
quantities are lar gely fixed once procurement and budgeting decisions are finalized.  
 
Bill of Quantities (BOQ) documents serve as foundational cost -planning instruments in construction 
projects. A BOQ provides itemized listings of materials, labor components, specifications, quantities, 
and associated costs. Despite containing rich quantita tive data, BOQs are typically treated as static 
financial documents rather than dynamic analytical datasets. Their tabular structure, when digitized 
and standardized, offers an opportunity to enable computational analysis and data -driven 
optimization.  
 
The rapid advancement of artificial intelligence (AI), machine learning (ML), and evolutionary 
optimization algorithms provides new possibilities for addressing complex, multi -criteria decision -
making problems in engineering. Construction planning inherent ly involves trade -offs between 
competing objectives, including cost, environmental impact, structural performance, and resource 
availability. Multi -objective optimization techniques allow simultaneous evaluation of conflicting 
goals and generate Pareto -optimal solution sets that provide decision -makers with flexible trade -off 
alternatives.  
 
Among these techniques, evolutionary algorithms such as the Non -dominated Sorting Genetic 
Algorithm II (NSGA -II) have demonstrated strong performance in solving non -linear, high -dimensional 
optimization problems. NSGA -II efficiently balances exploration an d exploitation through population -
based search mechanisms, non -dominated sorting, and crowding distance strategies. Its ability to 
generate a diverse Pareto frontier makes it particularly suitable for sustainable construction 
optimization, where no single solution can simultaneously minimi ze both cost and carbon without 
compromise.  
 
Recent research efforts have explored lifecycle assessment (LCA) integration within Building 
Information Modeling (BIM) systems, carbon estimation tools, and machine learning -based material 
prediction models. However, limited attention has been given to le veraging existing BOQ documents 
as primary analytical inputs for automated sustainability optimization. Furthermore, most existing 
studies focus either on carbon estimation or cost optimization independently, rather than integrating 
both within a unified A I-driven framework.  
 
This study addresses this research gap by proposing a scalable AI -based system that transforms raw 
BOQ data into structured datasets, maps materials to standardized emission factor databases, 
computes embodied carbon metrics, and applies multi -objective ev olutionary optimization to 
balance economic and environmental objectives. The framework aims to support early -stage 
construction planning decisions by providing actionable insights derived directly from conventional 
cost documentation.  
 
The primary contributions of this research are as follows:  
 
Development of an automated BOQ data structuring pipeline for sustainability analysis.  
 
Integration of emission factor mapping with cost modeling at item -level granularity.  
 
Implementation of NSGA -II for simultaneous cost and carbon optimization.  
 
Generation of Pareto -optimal trade -off solutions for informed stakeholder decision -making.  
 
Validation using realistic residential construction datasets.  
 
By embedding sustainability intelligence directly into cost -planning workflows, this approach 
contributes to the advancement of AI -enabled sustainable engineering practices. The framework 
supports regulatory compliance, carbon budgeting strategies, and lon g-term environmental 
responsibility without sacrificing economic feasibility.

\section{Methodology}
The proposed framework integrates data preprocessing, carbon computation, cost modeling, and
evolutionary multi -objective optimization into a unified pipeline. The overall architecture consists of 
four major stages: (1) BOQ data structuring, (2) emission f actor mapping and carbon calculation, (3) 
objective function modeling, and (4) multi -objective optimization using NSGA -II. Figure references 
are omitted here but can be added during formatting.  
 
1. BOQ Data Structuring and Preprocessing  
 
Bill of Quantities (BOQ) documents are typically available in spreadsheet or PDF formats and contain 
semi -structured tabular information such as material descriptions, quantities, units, and costs. 
However, these documents are not immediately suitable for computational analysis. Therefore, the 
first stage of the methodology focuses on transforming raw BOQ data into a clean, structured 
dataset.  
 
The preprocessing pipeline begins with document parsing. If the BOQ is in PDF format, text 
extraction tools are used to convert it into structured tabular data. Spreadsheet files are directly read 
and standardized. Each line item is converted into a consis tent data record containing the following 
attributes:  
 
Item identification number  
 
Material description  
 
Category (concrete, steel, masonry, finishing, etc.)  
 
Unit of measurement  
 
Quantity  
 
Unit rate  
 
Total cost  
 
Material descriptions often vary in wording. For example, “M25 grade concrete,” “reinforced cement 
concrete,” and “RCC mix” may refer to similar material categories. To address this, rule -based 
keyword matching and basic natural language processing techniq ues are applied to normalize 
descriptions and assign them to standardized material classes.  
Units are also standardized. For example, quantities expressed in tons are converted to kilograms, 
and cubic feet are converted to cubic meters where necessary. This ensures consistency during 
carbon and cost computation.  
The output of this stage is a clean dataset where every material entry is structured, categorized, and 
ready for environmental and economic analysis.  
 
2. Emission Factor Mapping and Carbon Computation  
 
Once the BOQ dataset is structured, each material is mapped to a standardized carbon emission 
factor obtained from lifecycle assessment databases. These emission factors represent the amount of 
carbon dioxide equivalent emitted per unit of material during extraction, production, and 
transportation.  
 
For example:  
Ordinary Portland Cement has a high emission factor due to energy -intensive clinker production.  
Recycled steel has a lower emission factor compared to virgin steel.  
Fly-ash bricks typically have lower embodied carbon than traditional clay bricks.  
After mapping emission factors to each material, embodied carbon is computed by multiplying the 
material quantity by its corresponding emission intensity. This calculation is performed at the item 
level and then aggregated to category level (e.g., total ca rbon from concrete works) and finally at the 
overall project level.  
Simultaneously, total project cost is calculated by aggregating the product of quantity and unit rate 
for each material.  
This dual computation enables the system to establish baseline values for:  
 
Total embodied carbon  
 
Total construction cost  
 
These baseline values serve as reference points for the optimization stage.  
 
3. Definition of Decision Variables  
 
To enable optimization, the system defines adjustable decision variables that influence cost and 
carbon outcomes. These variables include:  
 
Material substitution options  
High -carbon materials are assigned environmentally friendly alternatives. For example:  
 
Ordinary cement can be partially replaced with fly -ash blended cement.  
 
Conventional reinforcement steel can be replaced with recycled steel.  
 
Standard bricks can be replaced with fly -ash or compressed earth blocks.  
 
Quantity adjustment factors  
Minor structural or design optimizations can reduce excessive material usage within safety limits. The 
system allows small percentage -based quantity adjustments within predefined engineering 
tolerances.  
 
Constraint thresholds  
Cost variation limits are imposed to ensure economic feasibility. For instance, total project cost may 
not exceed a specified percentage above the baseline estimate. Additionally, minimum performance 
standards and material availability constraints are cons idered.  
 
These decision variables collectively form the solution space explored during optimization.  
 
4. Multi -Objective Optimization Using NSGA -II 
 
The core of the framework is the Non -dominated Sorting Genetic Algorithm II (NSGA -II), an 
evolutionary optimization algorithm designed for solving problems with multiple conflicting 
objectives.  
 
In this context, the two primary objectives are:  
 
Minimizing total construction cost  
 
Minimizing total embodied carbon emissions  
 
Unlike single -objective optimization, which produces one optimal solution, multi -objective 
optimization generates a set of trade -off solutions known as a Pareto frontier. Each solution on this 
frontier represents a configuration where improving one objecti ve would worsen the other.  
 
4.1 Population Initialization  
The algorithm begins by generating an initial population of candidate solutions. Each solution 
represents a unique combination of material substitutions and quantity adjustments.  
For example, one solution may substitute 30\% of cement with fly ash while keeping steel unchanged. 
Another solution may reduce concrete quantity slightly and replace steel entirely with recycled 
variants.  
 
4.2 Fitness Evaluation  
 
Each candidate solution is evaluated by recalculating:  
 
Total cost  
 
Total embodied carbon  
 
These values determine the solution’s fitness in the optimization process.  
 
4.3 Non -Dominated Sorting  
 
Solutions are ranked based on dominance relationships. A solution is considered superior if it 
performs better in at least one objective without performing worse in the other.  
 
The algorithm organizes solutions into different fronts:  
 
The first front contains non -dominated solutions (best trade -offs).  
 
Subsequent fronts contain progressively dominated solutions.  
 
4.4 Diversity Preservation  
 
To ensure a wide spread of solutions across the trade -off curve, NSGA -II uses a crowding distance 
mechanism. This prevents the algorithm from clustering around a narrow region of the solution 
space and encourages exploration of diverse design alternatives.  
 
4.5 Genetic Operators  
 
The algorithm evolves the population using:  
 
Selection mechanisms favoring higher -ranked solutions  
 
Crossover operations that combine features of two solutions  
 
Mutation operations introducing random variations  
 
This evolutionary process continues for multiple generations until convergence criteria are satisfied.  
 
5.Experimental Configuration  
 
The proposed framework was evaluated using real -world mid -scale residential construction BOQ 
datasets containing approximately 150 to 300 line items across structural, masonry, finishing, and 
reinforcement categories. The datasets were selected to reflect realistic material distributions, 
including high -carbon components such as cement -based concrete and structural steel. Before 
optimization, baseline project cost and embodied carbon were calculated to establish reference 
values for performance comparison.  
 
The evolutionary optimization process was configured with a population size of one hundred 
candidate solutions to ensure adequate exploration of the solution space while maintaining 
computational efficiency. The algorithm was executed for three hundred gen erations, which was 
determined sufficient for convergence based on preliminary experimental trials. A high crossover 
probability was used to promote effective combination of material substitution strategies between 
candidate solutions, while a moderate mutation rate was applied to preserve diversity and avoid 
premature convergence. Constraint limits were enforced during each generation to ensure that 
candidate solutions remained within acceptable cost deviation thres holds and engineering feasibility 
bounds.  
 
To evaluate performance, convergence behavior was monitored across generations by observing 
stability in Pareto -front distribution. Additional assessment metrics included percentage carbon 
reduction relative to baseline, percentage cost deviation, and dive rsity spread of solutions along the 
trade -off frontier. Computational runtime was also recorded to analyze scalability potential for larger 
BOQ datasets. The experimental setup demonstrated stable convergence and consistent generation 
of diverse trade -off configurations.  
 
6.Output Generation and Decision Support  
 
The final outcome of the optimization process is a Pareto frontier consisting of multiple non -
dominated solutions representing different trade -offs between total construction cost and embodied 
carbon emissions. Instead of producing a single optimal configu ration, the system generates a range 
of viable alternatives that reflect varying priorities between economic and environmental objectives.  
 
Each solution in the Pareto set includes a fully defined material configuration, updated quantity 
allocations, recalculated total project cost, and total embodied carbon value. These results are 
presented in a structured comparative format that enables sta keholders to directly evaluate the 
impact of each decision scenario. For instance, one solution may achieve moderate carbon reduction 
with negligible cost increase, while another may prioritize aggressive emission reduction within the 
upper boundary of the  allowed cost deviation.  
 
The decision -support capability of the framework lies in its ability to visualize trade -offs clearly and 
transparently. Project managers can assess whether marginal cost increases justify substantial carbon 
reductions. Policymakers can evaluate compliance with regulatory carbon targets. Engineers can 
analyze which material substitutions contribute most significantly to environmental improvement.  
 
By converting static BOQ documents into dynamic optimization outputs, the framework transforms 
traditional cost estimation into a strategic sustainability planning tool. The generated solutions 
empower decision -makers to balance financial constraints with long -term environmental 
responsibility, thereby promoting data -driven sustainable construction practices.

\section{Results}
Experimental evaluation was conducted using mid -scale residential BOQ datasets containing
between 150 and 300 line items. Baseline analysis indicated that concrete and steel components 
accounted for more than 60 percent of total embodied carbon. The NSGA -II optimization process 
generated a well -distributed Pareto frontier demonstrating clear trade -offs between cost and carbon 
reduction.  
 
Results show that embodied carbon reductions ranging from 18 percent to 32 percent were 
achievable while maintaining cost deviations within five percent of the baseline estimate. A balanced 
solution achieved approximately 24 percent carbon reduction with l ess than two percent increase in 
total project cost, indicating strong practical feasibility. Material substitutions, particularly partial 
cement replacement and recycled steel adoption, contributed significantly to emission reductions.  
 
The optimization framework demonstrated stable convergence and computational efficiency, 
confirming its applicability for real -world construction planning scenarios.

\section{Conclusion}
This study presented an AI -driven multi -objective optimization framework for sustainable
construction planning using structured BOQ datasets. By integrating emission factor mapping, 
embodied carbon computation, and evolutionary optimization through NSGA -II, the framework 
enables simultaneous consideration of economic and environmental objectives. Experimental results 
demonstrate that substantial carbon reductions can be achieved with minimal cost impact during 
early -stage project planning.  
 
The proposed approach transforms traditional BOQ documents from static cost estimation tools into 
intelligent sustainability decision -support systems. The generated Pareto -optimal solutions provide 
stakeholders with flexible trade -off options aligned with regulatory requirements and sustainability 
targets. Future research may focus on integration with Building Information Modeling platforms, 
real-time material pricing databases, and expanded lifecycle assessment parameters to enhance 
scalability and industr y adoption.

\section*{References}
\small
[1] Intergovernmental Panel on Climate Change, “Climate Change 2022: Mitigation of Climate
Change,” Cambridge University Press, 2022.  
 
[2] United Nations Environment Programme, “2023 Global Status Report for Buildings and 
Construction,” UNEP , 2023.  
 
[3] K. Deb, A. Pratap, S. Agarwal, and T. Meyarivan, “A fast and elitist multiobjective genetic 
algorithm: NSGA -II,” IEEE Transactions on Evolutionary Computation, vol. 6, no. 2, pp. 182 –197, 2002.  
 
[4] P . Cabeza, L. Rincón, V. Vilariño, G. Pérez, and A. Castell, “Life cycle assessment of building 
materials,” Building and Environment, vol. 44, no. 12, pp. 2510 –2516, 2009.

\end{document}
